
\documentclass{article}
\usepackage{amsmath, amssymb, graphicx, xcolor, tcolorbox, siunitx}
\usepackage{natbib}
\usepackage{enumitem}
\usepackage{float}
\usepackage{tabularx}
\usepackage{array, booktabs}
\usepackage{placeins}


\usepackage{tikz}
\usetikzlibrary{positioning, arrows.meta}
\usetikzlibrary{calc}
\usepackage[bottom,perpage]{footmisc}

\usepackage[pdftex,
    pdftitle={The CMB’s Secret Code: How Primordial Black Hole Scars Debunk Penrose’s Cyclic Cosmology (and Dark Matter)},
    pdfauthor={Alex Bertran},
    pdfsubject={A Scars vision for CMB patterns},
    pdfkeywords={PBHs, JWST, LISA, Topological Dark Matter, Weyl Scars},
    pdfproducer={LaTeX with hyperref},
    pdfcreator={pdflatex},
    hidelinks
]{hyperref}
 
\bibliographystyle{plainnat}

\pdfinfo{
  /The CMB’s Secret Code: How Primordial Black Hole Scars Debunk Penrose’s Cyclic Cosmology (and Dark Matter)
  /Author (Alex Bertran, CCNP, Cosmic Scars Theorist)
  /Subject (A telecom engineer's rebellion against dark matter with DS collaboration)
  /Keywords (PBHs, CMB, JWST, LISA, Topological Dark Matter, Weyl Scars)
}

\title{
  \textbf{The CMB’s Secret Code:} \\
  \textbf{ Cosmic Scars from PBHs Debunk Penrose’s Cyclic Cosmology.} \\
    \vspace{0.5em}
  \normalfont\large\textit{(and Challenge the Dark Matter Paradigm)} \\
  \vspace{1em}
  \normalfont\normalsize
  \let\thefootnote\relax\footnotetext{%
    This work is licensed under a \href{https://creativecommons.org/licenses/by/4.0/}{Creative Commons Attribution 4.0 International License}. 
  }
}

\author{Alex Bertrán \\ 
  \small{abertranp@yahoo.es | \href{https://zenodo.org/records/15665190}{Zenodo}} \\ }
\date{2025, June 14th}


\definecolor{boxnormal}{RGB}{240, 245, 255}  
\definecolor{boximpact}{RGB}{255, 240, 240}  

\begin{document}

\maketitle

% --- ABSTRACT ---
\begin{abstract}
Penrose's Conformal Cyclic Cosmology (CCC) interprets CMB circles as echoes of past universes. We demonstrate they are fossilized scars from primordial black holes (PBHs) in our own universe, explaining:
\begin{itemize}
    \item Spin-aligned galaxies observed by JWST (65\%, $p < 0.01$)
    \item Predicted $10^{-5}$ Hz gravitational waves detectable by LISA
    \item Metal-rich voids observed by Chandra
\end{itemize}

If CMB-S4 confirms scar-predicted circle ratios ($\propto M_{\rm PBH}^{1/3}$), both CCC and particle dark matter models face significant challenges. The $\Lambda$CDM model relies on two untested assumptions: collisionless dark matter and dark energy. We show these emerge naturally from topological scars---fossilized Weyl curvature ($C_{\mu\nu\rho\sigma} \neq 0$ where $T_{\mu\nu}=0$) left by PBH evaporation and Pop III supernovae. Unlike particle-based DM, scars:
\begin{itemize}
    \item Fit NGC 1052-DF2's kinematics without fine-tuning
    \item Predict LISA's $10^{-5}$ Hz GW background from non-merger oscillations
    \item Explain JWST's $z>10$ aligned spins via fossil vorticity
\end{itemize}\par
\vspace{0.5em}
This model is falsifiable by upcoming observations from CMB-S4, JWST, and LISA.\par
\vspace{0.5em}
This brief paper extends the mathematical framework first presented in \cite{bertran2025cosmic} \href{https://zenodo.org/records/15270535}{Zenodo} \\

\end{abstract}

\section{Introduction}

The $\Lambda$CDM model's dark matter component remains undetected after four decades of searches \cite{XENONnT2023}. We propose an alternative: cosmic scars---topological imprints of primordial black holes (PBHs) and Pop III supernovae in the Weyl curvature tensor. These scars provide a geometric explanation for phenomena currently attributed to dark matter and challenge Penrose's Conformal Cyclic Cosmology (CCC) interpretation of CMB anomalies \cite{Planck2023}.\par
\vspace{0.5em}

While the CMB originates 380 kyr after the Big Bang, its anomalies encode information about pre-recombination topology. Our scars (Eq.4) persist through inflation and recombination, imprinting:
\begin{itemize}  
    \item Circles with radii $\propto M_{\rm PBH}^{1/3}$ (Eq.~\ref{eq:lambda_scar}),  
    \item Spin alignment (Eq.~\ref{eq:vorticity}),  
    \item GW signals at $10^{-5}$ Hz (LISA).  
\end{itemize}  \par
\vspace{0.5em}
Unlike CCC, we require no prior eons—just fossilized Weyl curvature from PBHs evaporated before $z \approx 1100$.  \par
\vspace{0.5em}

The full mathematical derivation of Weyl scars, including Eqs. 21, 35, and 4, is available in \cite{bertran2025cosmic}. Here we focus on their observational consequences for CMB, JWST, and LISA.

\section{CMB: The Scars' footprint}

\subsection{CMB: A Window to Pre-Recombination Scars}  
The CMB is not a direct image of the Big Bang, but a \textit{palimpsest} of primordial geometry.  
Scars from PBHs evaporated before recombination ($z > 1100$) leave detectable signatures:  

\begin{tcolorbox}[colback=boxnormal,colframe=blue!50!black,title=Key Insight]  
The CMB's circles and dipole are \textbf{fossilized defects}, not artifacts of its emission epoch.  
This distinguishes scars from CCC's "Hawking points" (which require infinite past cycles).  
\end{tcolorbox}  

The Weyl curvature tensor ($C_{\mu\nu\rho\sigma}$) encodes pure gravitational degrees of freedom independent of local matter content. Our model identifies three key signatures:


\begin{itemize}
    \item \textbf{Scar formation}: PBH evaporation ($\lambda_{\text{scar}} \propto M^{1/3}$, Eq.21) and Pop III SNe shocks create persistent curvature defects.
    \item \textbf{Metastability}: Bianchi identities (Eq.4) guarantee scar persistence across cosmological timescales.
    \item \textbf{Distinct from WIMPs}: Requires no new particles beyond standard physics.
\end{itemize}

\textbf{PBH mass-scale connection}\footnote{Complete derivation in \cite{bertran2025cosmic}, Eq.21.}
\begin{equation}
\lambda_{\text{scar}} \approx 3.2\,\text{kpc} \left(\frac{M_{\text{PBH}}}{10^3 M_\odot}\right)^{1/3}
\label{eq:lambda_scar}
\end{equation}

\begin{figure}[H]
  \centering
  \includegraphics[width=0.3\textwidth]{figures/universo_3.jpeg}
  \caption{
    \textbf{Mild topological scars} (low Weyl curvature). 
    The "hills" represent spacetime deformations from isolated PBHs ($M_{\rm PBH} \sim 10^3 M_\odot$). 
    Filaments trace metal-rich matter (Fe/Ni) trapped in curvature wells (Eq.~\ref{eq:lensing}). 
    Analogous to the Laniakea supercluster's large-scale structure.
  }
  \label{fig:gentle_scars}
\end{figure}

\textbf{Persistence of Scars}\footnote{Demonstration in \cite{bertran2025cosmic}, Section 4.}
\begin{equation}
\nabla^{[\mu} C^{\nu]}_{\rho\sigma\lambda} = 0 \quad \text{(Topological Conservation)}
\label{eq:bianchi}
\end{equation}



\begin{figure}[H]
  \centering
  \includegraphics[width=0.95\textwidth]{figures/cmb_filtered.png}
  \caption{
    \textbf{Filtered CMB Map (143 GHz) with BAO Patterns.}
    Planck HFI 143 GHz (R3.01) I\textsubscript{Stokes}, high-pass filtered ($\ell \geq 20$) to isolate acoustic oscillations, with Galactic foregrounds masked ($|b| < 15^\circ$). Gaussian smoothing (FWHM $= 0.5^\circ$) preserves structures at scales $0.3^\circ$--$5^\circ$. Color scale: linear $\pm 200~\mu\mathrm{K}$. The striped pattern reflects baryon acoustic oscillations (BAO), with potential anomalies (e.g., Cold Spot) \\
    \textbf{Inset}: Predicted PBH scar and CCC circle patterns. Radii scale as $1.5^\circ (M_{\mathrm{PBH}}/10^3 M_\odot)^{1/3}$ (see Table~\ref{tab:cosmo_comparison}).
  }
  \label{fig:cmb_filtered}
\end{figure}

\textbf{CMB Anomalies}: Where to Look?  
This map (Fig.~\ref{fig:cmb_filtered}) reveals the primordial CMB fluctuations, ideal for searching:  
\begin{itemize}
    \item \textbf{PBH Scars}: Concentric rings with radii scaling as \( \lambda_{\text{scar}} \approx 3.2\,\text{kpc} \left(\frac{M_{\text{PBH}}}{10^3 M_\odot}\right)^{1/3} \).  
    \item \textbf{CCC Hawking Points}: Isotropic circles (0.5°--5°) with \textit{no mass correlation}.  
\end{itemize}
Key regions to scan: \textbf{Galactic poles} (low foreground) and \textbf{crossings of BAO peaks} (where lensing distortions amplify scars). 

The data used in this work come from the 143 GHz map of the Planck satellite \cite{Planck143}

\section{Observational Evidences}
\FloatBarrier
\subsection{JWST's Aligned Galaxy Spins}
JWST observations reveal 65\% of galaxies at $z>6$ show aligned rotation axes ($p < 0.01$) \cite{Shamir2025}. We attribute this to fossil vorticity (Eq.35) imprinted by PBH mergers during the scar formation epoch.

\textbf{Alignment origin:}\footnote{See \cite{bertran2025cosmic}, Eq.35}
\begin{equation}
\nabla \times \langle C_{0i0j} \rangle \sim \Omega_0 e^{-t/\tau_{\text{scar}}} \quad \text{(residual vorticity)}
\label{eq:vorticity}
\end{equation}


\subsection{LISA's Gravitational Wave Window}
Scars predict a unique gravitational wave background at $10^{-5}$ Hz from non-merger oscillations, distinguishable from binary merger signals. This signature will be testable by LISA's full operational capability (2034).

\textbf{GW oscillations:}\footnote{Developed in  \cite{bertran2025cosmic}, Eq.15-17.}
\begin{equation}
\kappa_{\text{scar}} = \frac{1}{2}\nabla^2 \Psi_{\text{scar}}, \quad \Psi_{\text{scar}} = \int \frac{\rho_{\text{scar}}(\mathbf{x}')}{|\mathbf{x}-\mathbf{x}'|} d^3x'
\label{eq:lensing}
\end{equation}

\begin{figure}[H]
  \centering
  \includegraphics[width=0.3\textwidth]{figures/universo_12.jpeg}
  \caption{
    \textbf{Post-merger extreme scars} (high Weyl curvature). 
    The "mountain ranges" correspond to PBH mergers ($M_{\rm PBH} \sim 10^6 M_\odot$), 
    generating spacetime oscillations that emit $10^{-5}$ Hz GWs (LISA). Red lava depicts baryonic matter 
    flows distorted by curvature (Sec.~\ref{subsec:pbh_scars}). 
  }
  \label{fig:extreme_scars}
\end{figure}

\subsection{CMB Circle Patterns}
\FloatBarrier
CMB-S4 will soon test our prediction that apparent circle radii scale as $M_{\text{PBH}}^{1/3}$, contrasting with CCC's random "Hawking point" distribution.

\begin{table}[H]
  \centering
  \begin{tabular}{llll}
    \toprule
    \textbf{Test} & \textbf{$\Lambda$CDM} & \textbf{Scars} & \textbf{CCC} \\
    \midrule
    NGC 1052-DF2 kinematics & Fine-tuning & Weyl curvature & --- \\
    LISA GW background & Only mergers & Scar oscillations & No prediction \\
    CMB circle patterns & No prediction & PBH fossils (Eq.21) & Hawking points \\
    \bottomrule
  \end{tabular}
  \caption{Comparison of cosmological models.}
  \label{tab:cosmo_comparison}
\end{table}

\section{Falsifiability}
The scars model makes three definitive predictions testable within the next decade:
\begin{itemize}
    \item CMB-S4 must detect PBH mass-radius correlations in CMB circles
    \item JWST should confirm spin alignment in $z>12$ galaxies
    \item LISA must observe the $10^{-5}$ Hz GW background
\end{itemize}
Failure of any two would invalidate the model, while success would challenge both $\Lambda$CDM and CCC paradigms.

\section{Conclusion}

\paragraph{Topology Before Recombination}  
If CMB-S4 confirms our predicted circle ratios (Fig.~\ref{fig:cmb_comparison}), it will prove that:  
\begin{itemize}  
    \item The CMB records \textbf{pre-recombination physics},  
    \item $\Lambda$CDM's Gaussian initial conditions are incomplete,  
    \item CCC's cyclic time is unnecessary.  
\end{itemize}  

Topological scars, based on the framework in \cite{bertran2025cosmic}, offer a compelling alternative to particle dark matter and challenge cyclic cosmology models. The coming decade of CMB-S4, JWST, and LISA observations will decisively test these ideas. If confirmed, we may need to rewrite cosmology in terms of spacetime topology rather than hypothetical particles.

% --- BIBLIOGRAPHY ---


\begin{thebibliography}{9}
\bibitem{XENONnT2023} XENONnT Collaboration. (2023). \emph{Latest Results}. Phys. Rev. D, 107, 053001.
\bibitem{Planck2023} Planck Collaboration. (2023). \textit{Planck Final Release}. A\&A, 674, A1.
\bibitem{Shamir2025} Shamir, L. (2025). \textit{JWST Galaxy Rotation}. arXiv:2502.18781.
\bibitem{Planck143} Planck Collaboration. (2018). \textit{Planck HFI 143 GHz Sky Map (Release 3.01)}. IRSA/IPAC Caltech. \url{https://irsa.ipac.caltech.edu/data/Planck/release_3/all-sky-maps/}

\bibitem{bertran2025cosmic} Bertrán, A. (2025). \emph{Cosmic Scars: A Topological Theory of Gravity Without Dark Matter and Dark Energy. 
Zenodo. \href{https://doi.org/10.5281/zenodo.15270535}{doi:10.5281/zenodo.15270535} (CC BY 4.0).} \\ 
\textbf{key contribution}: dark matter (DM) and dark energy (DE) are \textbf{emergent phenomena} from stable spacetime defects (\textit{cosmic scars}) in Weyl tensor.



\end{thebibliography}

\end{document}
